\documentclass[11pt,a4paper]{article}
\usepackage[margin=0.9in]{geometry}
\usepackage[T1]{fontenc}\usepackage[utf8]{inputenc}
\usepackage{xcolor}\usepackage[normalem]{ulem}\usepackage{amsmath,amssymb}
\usepackage{setspace}\usepackage{parskip}\usepackage{hyperref}
\hypersetup{colorlinks=true,linkcolor=blue,urlcolor=blue}
\definecolor{CutRed}{RGB}{200,30,30}
\definecolor{HeadBlue}{RGB}{0,70,180}
\definecolor{GrayNote}{RGB}{90,90,90}
\newcommand{\citep}[1]{\textsuperscript{[refs]}}
\newcommand{\figref}[1]{Fig.~X}
\newcommand{\cutx}[1]{{\color{CutRed}\sout{#1}}}
\newcommand{\newx}[1]{{\color{CutRed}\uline{#1}}}
\newcommand{\itemhead}[5]{\vspace{1.1em}\noindent{\color{HeadBlue}\large\textbf{#1}\quad\normalsize\textbf{#2}}\\[0.2em]
{\color{GrayNote}\small #3 words $\rightarrow$ #4 words, saves \textbf{#5}}\par}
\newcommand{\rationale}[1]{{\color{GrayNote}\small\textit{#1}}\par\smallskip}
\newcommand{\marknote}[1]{{\color{CutRed}\small\textbf{Reviewer-marked text involved:} #1}\par}
\newcommand{\panel}[1]{\par\smallskip\noindent\textbf{\small #1}\par}
\newcommand{\decision}{\par\smallskip\noindent{\color{GrayNote}\small$\square$~accept\quad$\square$~reject\quad$\square$~modify}\par
\noindent\rule{\linewidth}{0.4pt}}
\setstretch{1.05}
\begin{document}
\begin{center}{\Large\textbf{Tier 1 length reduction: side-by-side for author review}}\\[0.3em]
{\normalsize Prepared 2026-07-31. No manuscript file has been edited.}\end{center}
\vspace{0.5em}
{\small\textbf{How to read this.} For each item, \textbf{Current} shows the existing text with
\cutx{deleted wording struck through in red}; \textbf{Proposed} shows the resulting text with
\newx{newly written wording underlined in red}. Citations are collapsed to \citep{} and figure
references to \figref{} for readability; both are preserved verbatim in the real edit.
The red here is a review aid only and is unrelated to the manuscript's
\texttt{\textbackslash rev\{\}} / \texttt{\textbackslash revsecond\{\}} colour scheme.}
\vspace{0.3em}
\noindent\rule{\linewidth}{0.4pt}

\itemhead{I-1}{Introduction, closing paragraph}{197}{102}{95}

\rationale{Four preview sentences duplicate the Abstract (8 shared six-word sequences) and the Results subsection headings. The gap statement, the aims sentence, the first-use outcome definitions, and the closing reconciliation all stay.}
\panel{Current (197 words)}
Together, these contrasting lines of evidence indicate that communities behave as units of selection in some coalescence events but not others; however, the factors governing these transitions remain poorly understood. Here, we combine empirical and theoretical approaches to investigate how effective interaction intensity and environmental feedbacks drive these contrasting outcomes. Using randomly assembled synthetic microbial communities across different interaction regimes, we classify post-coalescence outcomes into three types: Dominance (one community wins), Mixture (both persist), and Restructuring (a novel state emerges). \cutx{Under weak interactions, Mixture prevails and species persistence is uncorrelated, consistent with individualistic dynamics. As effective interaction intensity increases, the system shifts toward Dominance, where species persistence becomes correlated, indicating community-level selection. A phenomenological generalized Lotka--Volterra (gLV) model with randomly sampled pairwise competitive interactions recapitulates the observed Dominance/Mixture transition, and a similar nutrient-dependent shift toward Dominance also appears in taxonomically richer communities derived from natural environmental samples. Further analysis reveals two mechanistic regimes underlying community-level selection: one where a few dominant taxa largely determine the winner, and another where collective multi-species dynamics shape the outcome. }These results reconcile conflicting observations by identifying effective interaction intensity, environmental feedbacks, and assembly-generated interaction structure as key determinants of community-level selection.
\panel{Proposed (102 words, saves 95)}
Together, these contrasting lines of evidence indicate that communities behave as units of selection in some coalescence events but not others; however, the factors governing these transitions remain poorly understood. Here, we combine empirical and theoretical approaches to investigate how effective interaction intensity and environmental feedbacks drive these contrasting outcomes. Using randomly assembled synthetic microbial communities across different interaction regimes, we classify post-coalescence outcomes into three types: Dominance (one community wins), Mixture (both persist), and Restructuring (a novel state emerges). These results reconcile conflicting observations by identifying effective interaction intensity, environmental feedbacks, and assembly-generated interaction structure as key determinants of community-level selection.
\decision

\itemhead{I-2}{Introduction, paragraph 2 (coalescence contexts)}{143}{113}{30}

\rationale{The four worked examples become a list. All seven citations are retained in place.}
\panel{Current (143 words)}
These contrasting frameworks make distinct predictions in the context of community coalescence, the mixing of previously isolated communities\citep{Rillig2015, Lechon2021}. \cutx{Coalescence occurs across diverse contexts and scales: environmental disturbances trigger microbial community mixing in soils\citep{Huet2023, Bresciani2025}, flooding events promote coalescence in aquatic and estuarine habitats\citep{Liu2024}, skin microbiomes undergo exchange through daily social contact\citep{Sarkar2024}, and gut microbiomes are subject to wholesale community transfer through fecal microbiota transplantation\citep{Xiao2020, Gupta2016}. }Coalescence brings species with distinct interaction histories into contact, creating novel cross-community interactions that can reshape the resulting assemblage\citep{Rillig2015}. Thus, it provides a natural test for the individualistic versus holistic paradigms: The holistic view predicts that species within a community have correlated persistence, making the community the primary unit of selection, termed community-level selection. The individualistic view predicts species-level selection, yielding outcomes shaped by individual fitness regardless of community origin\cutx{, with no systematic correlation in species persistence within parental communities}.
\panel{Proposed (113 words, saves 30)}
These contrasting frameworks make distinct predictions in the context of community coalescence, the mixing of previously isolated communities\citep{Rillig2015, Lechon2021}. \newx{Coalescence occurs across diverse contexts and scales, including soil disturbance\citep{Huet2023, Bresciani2025}, flooding of aquatic and estuarine habitats\citep{Liu2024}, social contact between skin microbiomes\citep{Sarkar2024}, and fecal microbiota transplantation\citep{Xiao2020, Gupta2016}. }Coalescence brings species with distinct interaction histories into contact, creating novel cross-community interactions that can reshape the resulting assemblage\citep{Rillig2015}. Thus, it provides a natural test for the individualistic versus holistic paradigms: The holistic view predicts that species within a community have correlated persistence, making the community the primary unit of selection, termed community-level selection. The individualistic view predicts species-level selection, yielding outcomes shaped by individual fitness regardless of community origin.
\decision

\itemhead{I-3}{Introduction, paragraph 1 (Clements and Gleason)}{107}{71}{36}

\rationale{Both paradigm names and all nine citations survive; the Discussion calls back to the Clements--Gleason debate. The closing sentence restates the opening question and is said again in the closing paragraph.}
\panel{Current (107 words)}
In nature, species coexist and interact within complex communities, yet whether these assemblages function as cohesive, integrated units or merely as loose collections of independently acting species remains a fundamental question in ecology. \cutx{Historically, this tension has been framed through two influential paradigms. }Clements' ``superorganism'' view treats communities as discrete biological entities with emergent properties arising from species interdependencies\cutx{, potentially as a result of coevolution}\citep{Clements1916, Clements1936, Odum1969, Lovelock1979, WilsonSober1989}. In contrast, Gleason's individualistic hypothesis posits that species independently occupy niches, with community composition emerging from the coincidental overlap of species ranges\citep{Gleason1939, Cain1947, Mason1947, Whittaker1967}. \cutx{Despite decades of research, the conditions that determine when communities behave as cohesive units versus loose species assemblages remain unclear.}
\panel{Proposed (71 words, saves 36)}
In nature, species coexist and interact within complex communities, yet whether these assemblages function as cohesive, integrated units or merely as loose collections of independently acting species remains a fundamental question in ecology. Clements' ``superorganism'' view treats communities as discrete biological entities with emergent properties arising from species interdependencies\citep{Clements1916, Clements1936, Odum1969, Lovelock1979, WilsonSober1989}. In contrast, Gleason's individualistic hypothesis posits that species independently occupy niches, with community composition emerging from the coincidental overlap of species ranges\citep{Gleason1939, Cain1947, Mason1947, Whittaker1967}. 
\decision

\itemhead{I-4}{Introduction, paragraph 3 (evidence for cohesion)}{150}{118}{32}

\rationale{The cut sentence's mechanism is stated again in Results section 2, where data support it. IMPORTANT: the Debray2022 citation must move onto the preceding sentence so no reference is lost.}
\panel{Current (150 words)}
Substantial theoretical and empirical evidence supports the holistic prediction that communities act as cohesive units during coalescence. \cutx{Early theoretical work by }Gilpin\citep{Gilpin1994} suggested that pre-assembled communities, having already undergone internal competitive exclusion, possess structured interaction patterns that differ fundamentally from randomly assembled communities. \cutx{Because species within such communities have been filtered to coexist, their survival outcomes become coupled, producing asymmetric post-coalescence communities dominated by one parental community\citep{Debray2022}. }Notably, this interaction structure can arise due to ecological exclusion alone, even in the absence of long-term coevolution, as demonstrated in synthetic microbial communities\citep{Venturelli2018, Maynard2020}. \cutx{Building on this foundation, resource-consumer} models formalized the mechanism by which such selection occurs, demonstrating that communities with superior collective fitness through resource consumption outcompete others in ways not predictable from individual species performance alone\citep{Tikhonov2016, Lechon2021, Xie2019, Xie2021}. Empirically, correlated selection between dominant and subdominant taxa has been observed across 100 coalescence experiments, confirming that species retention within communities is indeed correlated\citep{DiazColunga2022}.
\panel{Proposed (118 words, saves 32)}
Substantial theoretical and empirical evidence supports the holistic prediction that communities act as cohesive units during coalescence. \newx{}Gilpin\citep{Gilpin1994} suggested that pre-assembled communities, having already undergone internal competitive exclusion, possess structured interaction patterns that differ fundamentally from randomly assembled communities\newx{\citep{Debray2022}}. Notably, this interaction structure can arise due to ecological exclusion alone, even in the absence of long-term coevolution, as demonstrated in synthetic microbial communities\citep{Venturelli2018, Maynard2020}. \newx{Resource-consumer} models formalized the mechanism by which such selection occurs, demonstrating that communities with superior collective fitness through resource consumption outcompete others in ways not predictable from individual species performance alone\citep{Tikhonov2016, Lechon2021, Xie2019, Xie2021}. Empirically, correlated selection between dominant and subdominant taxa has been observed across 100 coalescence experiments, confirming that species retention within communities is indeed correlated\citep{DiazColunga2022}.
\decision

\itemhead{D-3}{Discussion, paragraph 2 closing sentence}{46}{15}{31}

\rationale{Restates paragraph 1's closing sentence about the Clements--Gleason debate. Verified not inside any \texttt{\textbackslash rev\{\}} span.}
\panel{Current (46 words)}
These differences may explain why community-level selection is more frequently observed in microbial coalescence studies. \cutx{This framework provides a unified perspective: rather than asking whether communities are cohesive units, we should ask under what conditions they become so, with interaction strength emerging as a key determinant.}
\panel{Proposed (15 words, saves 31)}
These differences may explain why community-level selection is more frequently observed in microbial coalescence studies. 
\decision

\itemhead{D-1}{Discussion, paragraph 3 (priority effects)}{71}{41}{30}

\rationale{Three sentences develop the parallel and then restate it. Both citations retained.}
\panel{Current (71 words)}
\cutx{The parallel with priority effects is particularly instructive; recent work showed that assembly order shapes final composition only under strong interactions, whereas weak interactions lead to convergent outcomes regardless of assembly history \citep{Hu2025, Fukami2015}. Our coalescence results mirror this pattern: strong interactions allow parental-community identity to shape the outcome, while weak interactions yield convergent Mixtures. Together, these findings reinforce interaction strength as a useful predictor of history-dependent community dynamics across multiple ecological contexts.}
\panel{Proposed (41 words, saves 30)}
\newx{The parallel with priority effects is instructive: assembly order shapes final composition only under strong interactions, whereas weak interactions give convergent outcomes regardless of assembly history \citep{Hu2025, Fukami2015}, mirroring our result that parental-community identity shapes the coalescence outcome only when interactions are strong.}
\decision

\itemhead{D-2}{Discussion, paragraph 5 (scope and outlook)}{132}{105}{27}
\marknote{first two sentences are \texttt{\textbackslash rev\{\}}}
\rationale{The first two sentences are reviewer-driven (\texttt{\textbackslash rev\{\}}, R2-1 relevant) and are untouched. The topic sentence is carried by the two sentences that follow it; the two habitat sentences merge.}
\panel{Current (132 words)}
Accordingly, we use community-level selection as an outcome-level description of origin-correlated persistence during coalescence. This operational definition does not by itself identify the causal ecological or biochemical mechanisms, which may vary across systems. \cutx{In nature, coalescence occurs in richer settings that may alter regimes and outcomes. Environmental heterogeneity (temperature, pH buffering, and other physicochemical factors) covaries with interaction strength and can reshape coalescence across habitats \citep{Rillig2015, Castledine2020, Rillig2017, Tropini2017}. In host-associated microbiomes, host filtering, priority effects, and spatial structure further influence successful colonization, making coalescence a useful framework for predicting and steering microbiome transfer \citep{Liu2024, Smillie2018, Rocca2021}. }Our experiments focus on steady-state outcomes and thus do not capture temporal dynamics such as fluctuating resources, migration pulses, or host responses. \cutx{Accounting for this broader ecological complexity will help build a more comprehensive picture of community-level selection across habitats and scales.}
\panel{Proposed (105 words, saves 27)}
Accordingly, we use community-level selection as an outcome-level description of origin-correlated persistence during coalescence. This operational definition does not by itself identify the causal ecological or biochemical mechanisms, which may vary across systems. \newx{In nature, coalescence occurs in richer settings: environmental heterogeneity covaries with interaction strength and can reshape coalescence across habitats \citep{Rillig2015, Castledine2020, Rillig2017, Tropini2017}, while in host-associated microbiomes host filtering, priority effects, and spatial structure further influence colonization \citep{Liu2024, Smillie2018, Rocca2021}. }Our experiments focus on steady-state outcomes and thus do not capture temporal dynamics such as fluctuating resources, migration pulses, or host responses. \newx{Accounting for this complexity will give a more comprehensive picture of community-level selection across habitats and scales.}
\decision

\itemhead{R-2}{Results section 1, paragraph 2 (experimental setup)}{232}{185}{47}

\rationale{The medium recipe and the phyla breakdown are reproduced word for word in Methods. All statistics, sample sizes, and figure references are kept.}
\panel{Current (232 words)}
To quantify the occurrence of outcome types experimentally, we assembled random \textit{in vitro} communities and subjected them to pairwise coalescence (\figref{fig:fig1}C). We first curated a strain library of 54 bacterial isolates collected from diverse environments (soil, tree surface, and flower stamen). \cutx{The library is phylogenetically broad, spanning 29 families across three phyla: Proteobacteria, Firmicutes, and Bacteroidota }(Extended Data Fig.~1; Supplementary Fig.~1). From this library, we assembled 30 parental communities with varying initial richness (6, 12, or 24 species; see Methods) and stabilized them for 7 days under daily growth--dilution cycles ($\times$30 every 24~h) (\figref{fig:fig1}C). Our initial experiments were done in Base medium \cutx{(1~g~L$^{-1}$ yeast extract, 1~g~L$^{-1}$ soytone, 10~mM sodium phosphate, 0.1~mM CaCl$_2$, 2~mM MgCl$_2$, 4~mg~L$^{-1}$ NiSO$_4$, 50~mg~L$^{-1}$ MnCl$_2$, 5~g~L$^{-1}$ glucose and 4~g~L$^{-1}$ urea; pH~6.5)}, a buffered complex medium that we have previously determined has moderate interaction strength and supports high species coexistence (Supplementary Information, mean species survival ratio = $74 \pm 2$\%). We then performed 83 pairwise coalescence events by mixing stabilized parental communities in equal volumes and restabilizing for an additional 7 days. Community compositions of parental and coalesced communities were measured at the end of stabilization by 16S rRNA amplicon sequencing (Amplicon Sequence Variants, ASVs), and we also recorded the optical density (OD$_{600}$) and pH of the communities to contextualize assembly and post-coalescence dynamics.
\panel{Proposed (185 words, saves 47)}
To quantify the occurrence of outcome types experimentally, we assembled random \textit{in vitro} communities and subjected them to pairwise coalescence (\figref{fig:fig1}C). We first curated a strain library of 54 bacterial isolates collected from diverse environments (soil, tree surface, and flower stamen). \newx{The library is phylogenetically broad }(Extended Data Fig.~1; Supplementary Fig.~1). From this library, we assembled 30 parental communities with varying initial richness (6, 12, or 24 species; see Methods) and stabilized them for 7 days under daily growth--dilution cycles ($\times$30 every 24~h) (\figref{fig:fig1}C). Our initial experiments were done in Base medium \newx{(Methods)}, a buffered complex medium that we have previously determined has moderate interaction strength and supports high species coexistence (Supplementary Information, mean species survival ratio = $74 \pm 2$\%). We then performed 83 pairwise coalescence events by mixing stabilized parental communities in equal volumes and restabilizing for an additional 7 days. Community compositions of parental and coalesced communities were measured at the end of stabilization by 16S rRNA amplicon sequencing (Amplicon Sequence Variants, ASVs), and we also recorded the optical density (OD$_{600}$) and pH of the communities to contextualize assembly and post-coalescence dynamics.
\decision

\itemhead{R-1}{Results section 1, paragraph 1 (framing and definitions)}{169}{129}{40}
\marknote{the Restructuring sentence is \texttt{\textbackslash rev\{\}}}
\rationale{The second Restructuring sentence restates the definition given one clause earlier. It is \texttt{\textbackslash rev\{\}} marked, so it is compressed rather than deleted. The R2-1 operational-definition sentence is untouched.}
\panel{Current (169 words)}
Our central question is whether selection acts primarily on whole communities as cohesive units during community coalescence. \cutx{To address this, we asked how often coalescence outcomes are represented primarily by one parental community versus blending of both or forming a novel community. }Consider a coalescence event where communities A and B merge to produce community C (\figref{fig:fig1}A). [\ldots equation and definitions unchanged \ldots] We partition this map into three coalescence outcome classes (see Methods): Dominance, where C closely resembles one parental community but not the other; Mixture, where C similarly resembles both parental communities; and Restructuring, where C diverges from both parental communities into a novel configuration. \cutx{Restructuring denotes low parental retention: the coalesced community is not well explained by the parental communities or their combination and may reflect a new stable composition produced by previously unseen cross-community species interactions. }Throughout, Dominance is the outcome-level signature that quantifies origin-correlated persistence from compositional structure, and community-level selection is our description of that persistence rather than a claim about a specific underlying mechanism.
\panel{Proposed (129 words, saves 40)}
Our central question is whether selection acts primarily on whole communities as cohesive units during community coalescence. Consider a coalescence event where communities A and B merge to produce community C (\figref{fig:fig1}A). [\ldots equation and definitions unchanged \ldots] We partition this map into three coalescence outcome classes (see Methods): Dominance, where C closely resembles one parental community but not the other; Mixture, where C similarly resembles both parental communities; and Restructuring, where C diverges from both parental communities into a novel configuration. \newx{Restructuring thus denotes low parental retention, possibly reflecting a new stable composition produced by previously unseen cross-community interactions. }Throughout, Dominance is the outcome-level signature that quantifies origin-correlated persistence from compositional structure, and community-level selection is our description of that persistence rather than a claim about a specific underlying mechanism.
\decision

\itemhead{R-3}{Results section 2, paragraph 1 (model setup)}{62}{41}{21}
\marknote{the caveat is \texttt{\textbackslash rev\{\}}}
\rationale{The scope caveat overlaps the Discussion's coarse-grained-descriptor sentence and is compressed, not deleted. The similarity-metric sentence restates Results section 1. The R1-7 four-community sentence is NOT touched (code audit pending).}
\panel{Current (62 words)}
[\ldots model equation, symbol definitions, and the four-community sentence unchanged \ldots] \cutx{Post-coalescence compositions were analyzed using the same similarity metrics as in the experiments. }\cutx{The gLV model serves as a phenomenological framework to explore how the statistical properties of interaction coefficients shape coalescence outcomes, rather than as a mechanistic model of the specific biochemical interactions (e.g., pH modification, metabolic cross-feeding) underlying our experimental observations.}
\panel{Proposed (41 words, saves 21)}
[\ldots model equation, symbol definitions, and the four-community sentence unchanged \ldots] \newx{Outcomes were analyzed with the same similarity metrics as the experiments. The gLV model is phenomenological, exploring how the statistical properties of interaction coefficients shape outcomes rather than modelling specific biochemical interactions.}
\decision

\vspace{0.5em}
\begin{center}\textbf{TOTALS}\end{center}
\noindent Tier 1 as drafted: \textbf{1309} words $\rightarrow$ \textbf{920} words, saving \textbf{389}.\par
\end{document}
